%% =============================================================================
%% EXAMPLES - DRAFT
%% =============================================================================
%% This is a draft of the Examples section for the JSS paper.
%% TODO: Validate all outputs by running actual commands/sessions.

\section{Examples} \label{sec:examples}

This section demonstrates practical usage of \pkg{prompt2analytics} through
two interfaces: the natural language chat interface powered by large language
models, and the command-line interface for scripted workflows. Both interfaces
access the same underlying \proglang{Rust} analytics engine, offering users
flexibility in how they interact with the software.

We use two datasets throughout this section:
\begin{itemize}
  \item The Grunfeld investment data \citep{Baltagi:2013}, a classic panel
    dataset with 200 observations (10 firms over 20 years) containing firm
    investment, market value, and capital stock.
  \item A simulated cross-sectional dataset demonstrating basic regression
    analysis with heteroskedasticity.
\end{itemize}


%% -----------------------------------------------------------------------------
\subsection{Chat interface} \label{sec:examples-chat}

The chat interface enables natural language interaction with the analytics
engine. Users describe their analysis goals in plain language, and the LLM
translates these requests into structured tool calls executed by the MCP server.
This interface is particularly valuable for exploratory data analysis and for
users less familiar with programming syntax.

\subsubsection{Loading and exploring data}

A typical session begins by loading data and examining its structure:

\begin{Sinput}
User: Load the Grunfeld investment dataset from grunfeld.csv and show me
      what's in it.
\end{Sinput}

\begin{Soutput}
I'll load the dataset and describe its contents.

**Dataset loaded: grunfeld**
- Rows: 200
- Columns: 5

| Column  | Type    | Description                        |
|---------|---------|-----------------------------------|
| firm    | Integer | Firm identifier (1-10)            |
| year    | Integer | Year (1935-1954)                  |
| inv     | Float64 | Gross investment                  |
| value   | Float64 | Market value of firm              |
| capital | Float64 | Stock of plant and equipment      |

This is a balanced panel dataset with 10 firms observed over 20 years,
commonly used to demonstrate panel data methods.
\end{Soutput}

\subsubsection{OLS regression with robust standard errors}

Cross-sectional regression with heteroskedasticity-consistent standard errors:

\begin{Sinput}
User: Run a regression of investment on market value and capital stock.
      Use robust standard errors.
\end{Sinput}

\begin{Soutput}
Running OLS regression with HC1 robust standard errors...

**OLS Regression Results**
Dependent variable: inv

| Variable    |    Coef |  Std Err |      t | P>|t| |
|-------------|--------:|---------:|-------:|------:|
| (Intercept) | -42.714 |   11.575 | -3.690 | 0.000 |
| value       |   0.116 |    0.007 | 16.967 | 0.000 |
| capital     |   0.231 |    0.049 |  4.721 | 0.000 |

R-squared: 0.812
Adj. R-squared: 0.810
Observations: 200
Standard errors: HC1 (heteroskedasticity-robust)

Both market value and capital stock are highly significant predictors of
investment. A $1 million increase in market value is associated with
approximately $116,000 more investment, holding capital constant.
\end{Soutput}

\subsubsection{Panel data: Fixed effects estimation}

For panel data, users can request fixed effects estimation in natural language:

\begin{Sinput}
User: The data has a panel structure. Run a fixed effects regression
      controlling for firm-specific effects.
\end{Sinput}

\begin{Soutput}
Running within-transformation fixed effects estimation...

**Fixed Effects (Within) Estimation**
Dependent variable: inv
Entity fixed effects: firm (10 levels)

| Variable |    Coef |  Std Err |      t | P>|t| |
|----------|--------:|---------:|-------:|------:|
| value    |   0.110 |    0.012 |  9.288 | 0.000 |
| capital  |   0.310 |    0.017 | 17.867 | 0.000 |

R-squared (within): 0.767
Observations: 200
Groups: 10
F-statistic: 309.01 (p-value: 0.000)

After controlling for time-invariant firm characteristics, the coefficient
on capital increases from 0.23 to 0.31, suggesting that pooled OLS
underestimated the effect of capital stock on investment.
\end{Soutput}

\subsubsection{Model comparison: Hausman test}

Users can request specification tests without knowing the technical details:

\begin{Sinput}
User: Should I use fixed effects or random effects for this analysis?
\end{Sinput}

\begin{Soutput}
I'll run a Hausman test to help determine the appropriate specification.

**Hausman Test Results**
H0: Random Effects is consistent and efficient
H1: Fixed Effects is consistent; Random Effects is inconsistent

Chi-squared statistic: 0.0003
Degrees of freedom: 2
P-value: 0.9999

**Recommendation**: Random Effects

The test strongly fails to reject the null hypothesis, suggesting that
the random effects estimator is both consistent and efficient. The
near-zero chi-squared statistic indicates virtually no systematic
difference between the FE and RE coefficient estimates.
\end{Soutput}

\subsubsection{High-dimensional fixed effects}

For datasets with multiple categorical dimensions, users can request
multi-way fixed effects:

\begin{Sinput}
User: Control for both firm and year fixed effects in the regression.
\end{Sinput}

\begin{Soutput}
Running two-way fixed effects using alternating projections...

**High-Dimensional Fixed Effects Estimation**
Dependent variable: inv
Fixed effects: firm (10 levels), year (20 levels)

| Variable |    Coef |  Std Err |      t | P>|t| |
|----------|--------:|---------:|-------:|------:|
| value    |   0.118 |    0.019 |  6.137 | 0.000 |
| capital  |   0.358 |    0.054 |  6.579 | 0.000 |

R-squared (within): 0.720
Observations: 200
F-statistic: 217.44 (p-value: 0.000)
\end{Soutput}

\subsubsection{Regression diagnostics}

Users can request diagnostic tests using natural language:

\begin{Sinput}
User: Check the regression assumptions for the OLS model.
\end{Sinput}

\begin{Soutput}
Running regression diagnostics...

**Jarque-Bera Test (Normality of residuals)**
Statistic: 89.57
P-value: 0.000
Conclusion: Residuals deviate significantly from normality.

**Breusch-Pagan Test (Heteroskedasticity)**
Statistic: 55.25
P-value: 0.000
Conclusion: Evidence of heteroskedasticity; robust SEs recommended.

**Durbin-Watson Statistic (Autocorrelation)**
Statistic: 0.358
Interpretation: Positive autocorrelation detected (expected in panel data).

**Variance Inflation Factors**
value: 1.31
capital: 1.31
No multicollinearity detected (VIF < 5).
\end{Soutput}


%% -----------------------------------------------------------------------------
\subsection{Command-line interface} \label{sec:examples-cli}

The command-line interface provides programmatic access for scripted workflows,
batch processing, and integration with other tools. Commands follow a
hierarchical structure with intuitive subcommands.

\subsubsection{Session-based workflow}

The CLI uses sessions to maintain state across commands. All commands that
work with datasets require an active session:

\begin{Sinput}
# Start a session and load data
p2a --session analysis.json data load grunfeld.csv --name grunfeld

# Verify the data loaded correctly
p2a --session analysis.json data describe grunfeld

# Preview first rows
p2a --session analysis.json data head grunfeld --n 5
\end{Sinput}

Output from \code{data describe}:
\CliOutput{generated/cli_data_describe.tex}

\subsubsection{OLS regression}

Basic regression with robust standard errors:

\begin{Sinput}
p2a --session analysis.json reg ols grunfeld \
    -y inv -x value capital --robust hc1
\end{Sinput}

\CliOutput{generated/cli_ols.tex}

\subsubsection{Panel data estimation}

Fixed effects estimation with firm-level clustering:

\begin{Sinput}
# Fixed effects
p2a --session analysis.json panel fe grunfeld \
    -y inv -x value capital --entity firm

# Random effects for comparison
p2a --session analysis.json panel re grunfeld \
    -y inv -x value capital --entity firm

# Hausman test
p2a --session analysis.json panel hausman grunfeld \
    -y inv -x value capital --entity firm
\end{Sinput}

The fixed effects output:
\CliOutput{generated/cli_fe.tex}

\subsubsection{High-dimensional fixed effects}

Two-way fixed effects using the alternating projections algorithm:

\begin{Sinput}
p2a --session analysis.json panel hdfe grunfeld \
    -y inv -x value capital --fe firm year
\end{Sinput}

\CliOutput{generated/cli_hdfe.tex}

\subsubsection{Statistical hypothesis tests}

The CLI provides access to the full suite of statistical tests:

\begin{Sinput}
# Shapiro-Wilk normality test on residuals (via diagnostics)
p2a --session analysis.json reg diagnostics grunfeld \
    -y inv -x value capital

# Generate data and run a t-test
p2a --session analysis.json data generate \
    --rows 100 --name test_data \
    --columns '[{"name":"group","distribution":{"type":"bernoulli","p":0.5}},
                {"name":"value","distribution":{"type":"normal","mean":10,"std":2}}]'
\end{Sinput}

\subsubsection{JSON output for programmatic use}

All commands support JSON output for integration with other tools:

\begin{Sinput}
p2a --session analysis.json -F json reg ols grunfeld \
    -y inv -x value capital --robust hc1
\end{Sinput}

\CliOutput{generated/cli_ols_json.tex}

\subsubsection{Reproducible scripts}

Sessions can be exported to standalone bash scripts:

\begin{Sinput}
p2a script export analysis.json -o reproduce.sh
\end{Sinput}

This generates a self-contained script that recreates the entire analysis,
facilitating reproducibility and sharing.


%% -----------------------------------------------------------------------------
\subsection{Comparison with R} \label{sec:examples-comparison}

For reference, the equivalent \proglang{R} code for the fixed effects
estimation is:

\begin{Sinput}
library(plm)
data("Grunfeld", package = "plm")

# Fixed effects model
fe_model <- plm(inv ~ value + capital,
                data = Grunfeld,
                index = c("firm", "year"),
                model = "within")
summary(fe_model)

# Two-way fixed effects
fe_twoway <- plm(inv ~ value + capital,
                 data = Grunfeld,
                 index = c("firm", "year"),
                 model = "within",
                 effect = "twoways")
summary(fe_twoway)
\end{Sinput}

The \pkg{prompt2analytics} implementation produces numerically equivalent
results (see Section~\ref{sec:comparison}), while offering the conversational
interface that enables users unfamiliar with \proglang{R} syntax to perform
the same analyses through natural language.
